    % Para eliminar la sangría
    [leftmargin=*]
    
    % Items con puntos
    \begin{itemize}[label=\textbullet]
        \item 
    \end{itemize}

    % Items con numeros con )
    \begin{enumerate}[label=\textbf{\arabic*)}] 
        \item
    \end{enumerate}

    % Items con letras misnusculas con )
    \begin{enumerate}[label=\textbf{\alph*)}] 
        \item
    \end{enumerate}

    % Imagen unica 
    \begin{figure}[H] \centering
        \includegraphics[width=0.5\linewidth]{}
        \caption{ } 
        \label{ }
    \end{figure}

    \begin{figure}
        \centering
        \resizebox{0.75\textwidth}{!}{\input{img/nrz}}
        \caption{}
        \label{}
    \end{figure}

    % Link
    \begin{center}
        \par \vspace{0.25 cm} \textbf{Para el desarrollo de esta experiencia se utilizó el siguiente} 
        \href{}
        {\textcolor{blue}{Google Colab.}}
    \vspace{0.25 cm}
    \end{center}
    
    % Doble imagen
    \begin{figure} [H] \centering
        \begin{subfigure} [H] {0.45 \linewidth} \centering
            \includegraphics[width=\linewidth]{}
            \caption{}
            \label{}
        \end{subfigure}
        \hspace{0.05 \linewidth}
        \begin{subfigure} [H] {0.45 \linewidth} \centering
            \includegraphics[width=\linewidth]{}
            \caption{}
            \label{}
        \end{subfigure}
        \caption{}
        \label{}
    \end{figure}

    % Imagen triple
    \begin{figure}[H]\centering
        \begin{subfigure}[b]{0.3\textwidth} 
            \includegraphics[width=\textwidth]{} 
            % \caption{}
            % \label{fig:}
        \end{subfigure}
        \hfill % Espacio horizontal entre imágenes
        \begin{subfigure}[b]{0.3\textwidth}
            \includegraphics[width=\textwidth]{} 
            % \caption{}
            % \label{fig:}
        \end{subfigure}
        \hfill
        \begin{subfigure}[b]{0.3\textwidth}
            \includegraphics[width=\textwidth]{} 
            % \caption{}
            % \label{fig:}
        \end{subfigure}
        % \caption{}
        % \label{}
    \end{figure}

    % Texto e imagen
    \begin{minipage}{0.5\textwidth}
        
    \end{minipage}
    \begin{minipage}{0.5\textwidth}
         \includegraphics[width=\linewidth]{}
    \end{minipage}

    % Tablas de las god
    \begin{table} [H] \centering
        \begin{tblr}{
            cells={ valign = m , halign = c },  % Alineación por defecto: izquierda
            colspec = {QQQQQ},                  % Define 5 columnas
            rowsep = 10 pt,                     % Espaciado entre filas
            colsep = 10 pt,                     % Espaciado entre columnas
            hlines,                             % Líneas horizontales
            vlines                              % Líneas verticales
        }

            % Aca va la tabla 
              
        \end{tblr}
        \caption{}  % El nombre de la tabla
        \label{}    % El tag de la tabla para la referencia
    \end{table}

    % idem pero una linea 
    \begin{table} [H] \centering
        \begin{tblr}{
            cells={ valign = m , halign = l }, colspec = {QQQ}, rowsep = 10 pt, colsep = 10 pt, hlines, vlines,
            column{1}= {bg=gray!15}, row{1} = {bold, halign=c} }

            % Aca va la tabla 
              
        \end{tblr}
        \caption{}  % El nombre de la tabla
        \label{}    % El tag de la tabla para la referencia
    \end{table}

    \begin{longtblr}[
        caption = titulo tabla,
        label = etiqueta
        ]{
            cells={ valign = m , halign = c }, 
            colspec = {Q Q[11 cm] Q}, 
            rowsep = 5 pt,        
            colsep = 7.5 pt,           
            hlines,              
            vlines,                             
            row{1}={bg=gray!15, font=\bfseries}
        }

            aa & aa & aa \\

    \end{longtblr}
    